% generated by sdr.manuscript; do not edit
{\footnotesize
\begin{longtable}{>{\raggedright\arraybackslash}p{\dimexpr0.2400\linewidth-2\tabcolsep\relax}>{\raggedright\arraybackslash}p{\dimexpr0.4600\linewidth-2\tabcolsep\relax}>{\raggedright\arraybackslash}p{\dimexpr0.3000\linewidth-2\tabcolsep\relax}}
\caption{Data items of the charting form (charting\_schema.json).}\label{tab:items}\\
\toprule
\textbf{Field} & \textbf{Definition} & \textbf{Allowed values} \\ \midrule
\endfirsthead
\multicolumn{3}{l}{\emph{\tablename~\thetable{} (continued)}}\\
\toprule
\textbf{Field} & \textbf{Definition} & \textbf{Allowed values} \\ \midrule
\endhead
\midrule\multicolumn{3}{r}{\emph{continued}}\\
\endfoot
\bottomrule
\endlastfoot
\texttt{id} & Stable slug assigned at record-to-dataset mapping (protocol 7.\allowbreak{}6), e.\allowbreak{}g. & dataset\_\allowbreak{}id \\
\texttt{name} & Name as the primary describing record gives it. & string \\
\texttt{aliases} & Other names, challenge names, repository titles. & array \\
\texttt{versions} & Named versions when any exist (VerSe 2019, VerSe 2020; TotalSegmentator v1, v2). & array \\
\texttt{charted\_\allowbreak{}version} &  & string \\
\texttt{modality} & Every modality the dataset contains. & CT / CBCT / radiograph / fluoroscopy / DXA-VFA / EOS / MRI / other \\
\texttt{projection} & For projection radiography. & lateral / AP / PA / oblique / biplanar / whole-spine / other / NA / unknown \\
\texttt{region} & Multi-select coverage: list every region (cervical, thoracic, lumbar, sacrum-pelvis) the images cover in every subject. & cervical / thoracic / lumbar / sacrum-pelvis / whole-spine / mixed / unknown \\
\texttt{n\_\allowbreak{}subjects} & Distinct persons in the eligible modalities. & integer | null \\
\texttt{n\_\allowbreak{}subjects\_\allowbreak{}basis} & 'computed-from-metadata' = counted from released patient identifiers; 'estimated' = derived from a stated ratio (e.\allowbreak{}g. & reported / computed-from-metadata / estimated / unknown \\
\texttt{n\_\allowbreak{}images\_\allowbreak{}or\_\allowbreak{}volumes} &  & integer | null \\
\texttt{n\_\allowbreak{}images\_\allowbreak{}unit} &  & images / volumes / series / studies / figures / sub-figures / unknown \\
\texttt{n\_\allowbreak{}per\_\allowbreak{}modality} & Subject and image counts per eligible modality when the dataset is multi-modal; keys are modality enum values. & object \\
\texttt{n\_\allowbreak{}annotated\_\allowbreak{}structures} & Total annotated structures (vertebrae, landmarks) in the release when stated; the denominator for annotation\_\allowbreak{}validation.\allowbreak{}fraction\_\allowbreak{}of\_\allowbreak{}items when the statistic is per structure. & integer | null \\
\texttt{n\_\allowbreak{}use\_\allowbreak{}records} & Records coded use-only that map to this dataset (charting/\allowbreak{}dataset\_\allowbreak{}sources.\allowbreak{}csv role 'use'); filled by code. & integer \\
\texttt{subject\_\allowbreak{}overlap} & Known or suspected sharing of subjects with another included dataset that is not a parent/\allowbreak{}derived relation (shared source CTs, reused NHANES films). & array \\
\texttt{pediatric\_\allowbreak{}included} & 'yes' when any subject under 18 years is stated to be included. & yes\_\allowbreak{}no\_\allowbreak{}unknown \\
\texttt{age\_\allowbreak{}range} &  & string \\
\texttt{annotation\_\allowbreak{}types} & 'none' must be the only element when present. & annotation\_\allowbreak{}type \\
\texttt{annotation\_\allowbreak{}origin} & Union of origins over all annotation types (protocol 4.\allowbreak{}3 definitions). & origin \\
\texttt{annotation\_\allowbreak{}origin\_\allowbreak{}by\_\allowbreak{}type} & Origin per annotation type, keys from the annotation\_\allowbreak{}type enum. & object \\
\texttt{annotator\_\allowbreak{}qualification\_\allowbreak{}stated} & True when the description names the annotators' role or qualification (radiologist, spine surgeon, 'trained annotator supervised by a radiologist') for the annotation type the rubric scores. & boolean \\
\texttt{annotation\_\allowbreak{}validation} &  & object: described, qa\_\allowbreak{}process, inter\_\allowbreak{}rater\_\allowbreak{}statistic, inter\_\allowbreak{}rater\_\allowbreak{}value, inter\_\allowbreak{}rater\_\allowbreak{}ci, inter\_\allowbreak{}rater\_\allowbreak{}sd, inter\_\allowbreak{}rater\_\allowbreak{}n, fraction\_\allowbreak{}of\_\allowbreak{}items, per\_\allowbreak{}structure\_\allowbreak{}breakdown, n\_\allowbreak{}raters, comparator, n\_\allowbreak{}categories, comparable\_\allowbreak{}group \\
\texttt{variant\_\allowbreak{}labels} &  & object: L6, T13, LSTV, other \\
\texttt{per\_\allowbreak{}vertebra\_\allowbreak{}level\_\allowbreak{}names} & 'yes' when every annotated vertebra carries an anatomical level name (C1.\allowbreak{}.\allowbreak{}S1) in the release; 'partial' when only some images or only one region do; 'no' when labels are class-agnostic or absent. & yes / partial / no \\
\texttt{documentation} & Read by rubric item Q1. & object: article\_\allowbreak{}or\_\allowbreak{}preprint, repository\_\allowbreak{}record\_\allowbreak{}or\_\allowbreak{}readme, description\_\allowbreak{}at\_\allowbreak{}least\_\allowbreak{}100\_\allowbreak{}words, states\_\allowbreak{}population, states\_\allowbreak{}acquisition, states\_\allowbreak{}annotation\_\allowbreak{}protocol, states\_\allowbreak{}file\_\allowbreak{}layout\_\allowbreak{}and\_\allowbreak{}label\_\allowbreak{}scheme \\
\texttt{metadata\_\allowbreak{}elements} & Read by rubric item Q5. & age-sex / scanner-or-acquisition-parameters / spacing-or-dicom-headers / site-and-country / inclusion-exclusion-criteria / clinical-status-per-subject \\
\texttt{licence} &  & object: text, tier, url, spdx, standard, stated\_\allowbreak{}in, consistent \\
\texttt{described\_\allowbreak{}as\_\allowbreak{}public} & Reviewer-charted, quote-required (protocol 4.\allowbreak{}3): the record states that the images are released, or gives a repository DOI or URL for the images (not for code or supplements). & yes\_\allowbreak{}no\_\allowbreak{}unknown \\
\texttt{described\_\allowbreak{}as\_\allowbreak{}public\_\allowbreak{}regex} & Filled by code: PUBLIC\_\allowbreak{}RE (protocol 4.\allowbreak{}3) matches title + abstract + description. & boolean \\
\texttt{access\_\allowbreak{}url} & The landing URL probed: the URL the description gives for the images (the first that is not the article itself), or the URL resolved as access\_\allowbreak{}url\_\allowbreak{}basis says. & string | null \\
\texttt{access\_\allowbreak{}url\_\allowbreak{}basis} & How access\_\allowbreak{}url was obtained (protocol section 10). & given-in-record / resolved-from-named-repository / resolved-from-doi / none \\
\texttt{host\_\allowbreak{}type} & Filled by code from the access\_\allowbreak{}url host through the protocol section-10 host table; the stratifier for SAP A3.\allowbreak{}1. & repository-api / institutional / personal-or-cloud-drive / challenge-platform / unknown \\
\texttt{file\_\allowbreak{}url\_\allowbreak{}probed} & One file URL from enrichment, probed with a ranged GET. & string | null \\
\texttt{access\_\allowbreak{}probe} & One entry per probe run. & probe \\
\texttt{access\_\allowbreak{}route} & Final route: the machine classification the two probes agree on, or the verifier's decision (protocol section 10). & access\_\allowbreak{}route \\
\texttt{access\_\allowbreak{}route\_\allowbreak{}rule} & The rule ids of the agreeing probes, slash-joined when the two probes reached the same classification by different rules (e.\allowbreak{}g. & string \\
\texttt{year} & Year of the earliest record that describes the image release (protocol 4.\allowbreak{}1, I5), not of the cohort. & integer \\
\texttt{year\_\allowbreak{}basis} &  & article / preprint / repository-record / challenge-page / unknown \\
\texttt{cohort\_\allowbreak{}year\_\allowbreak{}range} & Years in which the images were acquired, as stated (e.\allowbreak{}g. & string | null \\
\texttt{country} & ISO 3166-1 alpha-2 of the acquiring institution(s); more than one entry when multi-national. & array \\
\texttt{source\_\allowbreak{}record} &  & object: doi, pmid, pmcid, arxiv, url, title, record\_\allowbreak{}type, record\_\allowbreak{}ids \\
\texttt{size\_\allowbreak{}gb} & Total download size as stated or as summed from the repository API; null when unknown. & number | null \\
\texttt{format} &  & DICOM / NIfTI / NRRD / MHA / PNG / JPEG / TIFF / BMP / NPZ / HDF5 / MAT / parquet / raw / other / unknown \\
\texttt{versioning} & Read by rubric item Q8. & object: dataset\_\allowbreak{}doi, doi\_\allowbreak{}resolves\_\allowbreak{}to, version\_\allowbreak{}identifier, changelog\_\allowbreak{}or\_\allowbreak{}versioned\_\allowbreak{}records \\
\texttt{superset\_\allowbreak{}of\_\allowbreak{}version} & When the charted (largest) version is stated to contain an earlier release, that release's label (protocol 7.\allowbreak{}6); null otherwise. & string | null \\
\texttt{split\_\allowbreak{}defined} & A train/\allowbreak{}validation/\allowbreak{}test partition is released or specified. & yes\_\allowbreak{}no\_\allowbreak{}unknown \\
\texttt{patient\_\allowbreak{}level\_\allowbreak{}dedup} & The split (or the dataset) is stated to keep every subject in one partition and to contain each subject once. & yes / no / unknown / NA \\
\texttt{grouping\_\allowbreak{}variable\_\allowbreak{}released} & A patient id or equivalent is released so the grouping can be checked. & boolean \\
\texttt{image\_\allowbreak{}quality\_\allowbreak{}reported} & quantitative = per-image or per-series metrics, reader grades or artefact flags released; qualitative = a statement of quality-based exclusions with counts or a qualitative description. & yes-quantitative / yes-qualitative / no / unknown \\
\texttt{resolution\_\allowbreak{}summary} & Matrix size, pixel spacing, slice thickness, kV/\allowbreak{}mAs as reported; null when none. & string | null \\
\texttt{flags} &  & object: figure\_\allowbreak{}derived, maintainer\_\allowbreak{}authors, derived, parent\_\allowbreak{}ids, multi\_\allowbreak{}modal, challenge, authors\_\allowbreak{}own, superseded \\
\texttt{checksums} & Published checksums or sizes pulled by sdr/\allowbreak{}enrich.\allowbreak{}py; used by the recipe for verification. & checksum \\
\texttt{on\_\allowbreak{}grid} &  & object: status, path, checked\_\allowbreak{}date, n\_\allowbreak{}subjects\_\allowbreak{}on\_\allowbreak{}grid \\
\texttt{field\_\allowbreak{}sources} & Field name -> where its value came from: 'api:<source>', 'article', 'repository-page', 'reviewer', or 'disk:<path>@<sha1>' for values read from files on the grid. & object \\
\texttt{notes} &  & string \\
\texttt{quotes} &  & quote \\
\texttt{charted\_\allowbreak{}by} &  & string \\
\texttt{charted\_\allowbreak{}date} &  & iso\_\allowbreak{}date \\
\texttt{verified\_\allowbreak{}by} & Null until verified; verification is required for every dataset before analysis (protocol section 8) and, for authors'-own datasets, must not be in flags.\allowbreak{}maintainer\_\allowbreak{}authors. & string | null \\
\texttt{verified\_\allowbreak{}date} &  &  \\
\texttt{eligibility} &  & object: decision, criteria\_\allowbreak{}met, found\_\allowbreak{}via \\
\end{longtable}}
